\documentclass[../main.tex]{subfiles}

\begin{document}

%% ======================================================================
%%  Section 9 – Long-term Data Accessibility and Interpretability
%% ======================================================================
\section{Long-term Data Accessibility and Interpretability}
\label{sec:long-haul}

Long-term accessibility depends on more than preservation of the primary file
bytes. It can also depend on format bounds, required filters and extensions,
external resources, application schema and provenance, and the continued
availability of readers that implement the relevant interpretation. Community
discussion illustrates the operational importance of preserving plugin and
reader knowledge with archived data~\cite{forum-plugin}. This assessment did
not inventory an archive or test a representative preservation horizon.

\subsection{Audit Findings and Mitigation Planning}

The present chapter does not yet define an affected archive, time horizon,
failure condition, impact, or exposure sufficient for a separate formal
finding. The bounded reader/writer compatibility cases that bear on long-term
access are registered as \code{IND-SAF-001} and \code{IND-SAF-002} in
Appendix~\ref{app:findings-register}.

\paragraph{Priority reconciliation.}\mbox{}\\
This chapter therefore creates no additional response route. Treatment of the
two bounded compatibility cases follows their Planned clock and the acceptance
criteria in \code{RM-2C}; broader archival concerns require a defined archive,
failure condition, consequence, and exposure before they can be rated or
prioritized.

\end{document}
